\section{Compliance}
\label{sec:compliance}

\legal{This section describes a capability that is built but has not yet had
legal review. The mechanism is decided (\dref{D6}); the policy under which it
will be used, the entity that will operate it and its obligations are
subject to legal review before the chain ships. Nothing here is a
representation about the regulatory status of the network or its token.}

\subsection{What the chain can do}

Konstellation includes a chain-level compliance module, \code{x/compliance},
holding two lists (\dref{D6}, ENGINEERING \S10):

\begin{description}[style=nextline]
\item[An allowlist (\code{isVerified}).] Addresses that an attester has marked
  as verified. \textbf{The chain never requires allowlisting to transact.} The
  list exists so that a contract that wants to serve only verified addresses
  --- a permissioned pool, a regulated asset --- can consult it. Its use is
  entirely the contract author's choice.
\item[A blocklist (\code{isFrozen}).] Addresses the chain refuses to transact
  with. This list is enforced \textbf{chain-wide}: every transaction touching
  a listed address --- EVM or Cosmos, native \kash{} included --- is refused
  before execution, at the ante handler, and a synchronous mempool pre-check
  returns ``address is frozen'' at the JSON-RPC so the sender learns
  immediately. An incoming IBC transfer to a frozen address is refused at
  receipt and refunded by the sending chain.
\end{description}

Both lists are readable from Solidity through a read-only precompile,
\code{ICompliance}, at the fixed address
\code{0x0000000000000000000000000000000000000900}, exposing
\code{isVerified(address)}, \code{isFrozen(address)} and
\code{status(address)}. Contracts that need to refuse \emph{internal}
transfers to frozen addresses --- which the ante handler cannot see --- call
it.

One EVM-specific detail matters for wallets: adding an address to the
blocklist also clears any EIP-7702 delegation on it. Without that, an
externally owned account delegated to a smart-account wallet before the freeze
would still run that wallet's code when an EntryPoint or forwarder called it,
and a relayer could drain it with an operation the frozen key had signed
off-chain. Only the 23-byte delegation designator is cleared --- never real
contract bytecode, so freezing a token contract does not brick its holders ---
and storage is untouched, so a lifted account is restored by a single new
authorisation (ENGINEERING \S10, \S\ref{sec:aa}).

\subsection{Who can change the lists, and how fast}
\label{sec:complianceauth}

\begin{table}[H]
\centering
\small
\begin{tabularx}{\textwidth}{@{}lLL@{}}
\toprule
Path & Who & Timing \\
\midrule
Scheduled change & the compliance authority (a foundation multisig) & takes effect after an on-chain \textbf{timelock of 24 hours} on mainnet; visible on-chain for the whole delay \\
Emergency freeze & the compliance authority & takes effect \textbf{immediately}; \textbf{auto-expires after one timelock period} unless ratified by scheduling the permanent add. An address under a live emergency freeze, or within one timelock after it lapsed or was lifted, cannot be emergency-frozen again --- only the scheduled path or governance can hold a freeze past one period \\
Governance override & any passed governance proposal & can add or remove any entry, cancel any scheduled change, and \textbf{replace the authority} \\
\bottomrule
\end{tabularx}
\caption{Compliance list authority (\dref{D6}). On \code{testnet-1} the authority
is a development key with short timelocks; on mainnet it is the foundation
multisig with 24-hour timelocks, set in the network's genesis after legal
review (ENGINEERING \S18).}
\end{table}

Every action on either list is an on-chain event. The timelock is bounded in
code (one minute to ninety days) so that no parameter change can make the
lists unreachable, and the ratification path is designed so a freeze never
lapses in the gap between an emergency action and its permanent form. If no
authority is set, only governance can touch the lists.

\todo{The mainnet compliance authority multisig --- its signers, threshold and
the legal entity behind it --- is not yet chosen. It is a genesis entry.}

\subsection{What the chain deliberately does not do}

\begin{itemize}
\item It does not \textbf{reverse} transactions or move balances. There is no
  protocol-level clawback. The design options for reversal (issuer mint/burn of
  controlled assets, delayed settlement above a threshold, a per-token clawback
  role, or a governance state migration) were catalogued during design
  (ENGINEERING \S10); none is part of this release, and any asset-level control
  would be the issuer's contract, not the chain's.
\item It does not require identity to use the chain. The allowlist is opt-in for
  contracts; it gates nothing at the protocol level.
\item It does not store personal data. \todo{The engineering record describes a
  design in which attestations (credential hash, issuer, level, expiry,
  jurisdiction) are held on-chain and PII off-chain by a licensed provider. Whether
  the allowlist will carry attestation metadata beyond a boolean, and who the
  attesters are, is not decided.}
\end{itemize}

\subsection{The trade-offs, stated in full}
\label{sec:compliancerisks}

The engineering record requires that the following risks be recorded in this
document if the chain-wide design was chosen. It was, and they apply without
qualification (ENGINEERING \S10, \dref{D6}).

\begin{enumerate}
\item \textbf{The freeze authority becomes the highest-value key on the chain.}
  It is a multisig from genesis; non-emergency actions are timelocked; every
  action is logged on-chain; governance can replace it. Those mitigations
  reduce the risk; they do not remove it. A compromise of the authority could
  freeze any address for one timelock period before governance could act.
\item \legal{\textbf{The ability to freeze can create a legal obligation to
  freeze}, across multiple jurisdictions. Operating the authority is an
  ongoing, staffed function with liability attached, not a one-time deployment.
  Who bears that obligation, under which laws, and with what process for
  requests from authorities, has not been determined.}
\item \textbf{It permanently changes the chain's positioning.} Some developers
  will not build on a chain with a freeze switch; some institutions will not
  touch one without. Konstellation has chosen the second audience knowingly.
  Readers who need a chain with no such control should treat this as
  disqualifying rather than expect it to be removed.
\end{enumerate}

\legal{\todo{Legal review of \S\ref{sec:compliance} as a whole is required
before this document is released and before the module ships on mainnet:
the operating policy for the lists, the request-handling process, the
disclosure language above, and the regulatory implications of the allowlist
and any attestation scheme.}}
